%========================================
\chapter{附录：常见考点与公式汇总}
%========================================

\section{重要公式汇总}

\subsection{物理层}

\begin{formulabox}
\textbf{奈氏准则}（无噪声信道）：
$$B_{max} = 2W \text{ Baud}$$
其中$W$为带宽(Hz)

\textbf{香农公式}（有噪声信道）：
$$C = W \log_2\left(1 + \frac{S}{N}\right) \text{ bps}$$

\textbf{信噪比转换}：
$$\text{信噪比(dB)} = 10\log_{10}\left(\frac{S}{N}\right)$$

\textbf{比特率与波特率}：
$$\text{比特率} = \text{波特率} \times \log_2(\text{码元种类数})$$

\textbf{传输时延}：
$$T_s = \frac{\text{数据量(bit)}}{\text{发送速率(bps)}}$$

\textbf{传播时延}：
$$T_p = \frac{\text{距离(m)}}{\text{传播速度(m/s)}}$$

\textbf{总时延}：
$$T = T_s + T_p + T_{处理} + T_{排队}$$

\textbf{时延带宽积}：
$$\text{时延带宽积} = T_p \times \text{带宽}$$
\end{formulabox}

\subsection{数据链路层}

\begin{formulabox}
\textbf{停等协议信道利用率}：
$$U = \frac{T_D}{T_D + RTT}$$

\textbf{CSMA/CD争用期}：
$$\text{争用期} = 2\tau = \frac{2d}{v}$$

\textbf{最小帧长}：
$$L_{min} = 2\tau \times R$$

\textbf{GBN发送窗口}：
$$1 < W_T \leq 2^n - 1$$

\textbf{SR发送窗口}：
$$1 < W_T \leq 2^{n-1}$$
\end{formulabox}

\subsection{网络层}

\begin{formulabox}
\textbf{子网主机数}：
$$\text{主机数} = 2^{\text{主机位数}} - 2$$

\textbf{子网数}：
$$\text{子网数} = 2^{\text{子网位数}}$$

\textbf{IP分片片偏移}：
$$\text{片偏移} = \frac{\text{数据偏移量}}{8}$$
\end{formulabox}

\subsection{传输层}

\begin{formulabox}
\textbf{TCP序号}：
$$\text{下一序号} = \text{当前序号} + \text{数据字节数}$$

\textbf{TCP确认号}：
$$\text{确认号} = \text{期望收到的下一字节序号}$$

\textbf{TCP超时重传时间}：
$$\text{RTO} = \text{RTT}_S + 4 \times \text{RTT}_D$$

\textbf{TCP拥塞窗口}：
慢开始：$cwnd = cwnd \times 2$（指数增长）

拥塞避免：$cwnd = cwnd + 1$（线性增长）
\end{formulabox}

\subsection{应用层}

\begin{formulabox}
\textbf{HTTP传输时间}：
$$T = \text{RTT} \times \text{往返次数} + \text{传输时延}$$

\textbf{DNS查询时间}：
递归查询：$T = n \times \text{RTT}$（n为层级数）

迭代查询：$T = (2n-1) \times \text{RTT}$
\end{formulabox}

\section{常见端口号汇总}

\begin{table}[htbp]
\centering
\begin{tabular}{llll}
\toprule
\textbf{协议} & \textbf{端口} & \textbf{传输层} & \textbf{功能} \\
\midrule
FTP数据 & 20 & TCP & 文件传输数据连接 \\
FTP控制 & 21 & TCP & 文件传输控制连接 \\
SSH & 22 & TCP & 安全远程登录 \\
Telnet & 23 & TCP & 远程登录 \\
SMTP & 25 & TCP & 发送邮件 \\
DNS & 53 & UDP/TCP & 域名解析 \\
DHCP服务器 & 67 & UDP & 动态主机配置 \\
DHCP客户端 & 68 & UDP & 动态主机配置 \\
TFTP & 69 & UDP & 简单文件传输 \\
HTTP & 80 & TCP & Web服务 \\
POP3 & 110 & TCP & 接收邮件 \\
SNMP & 161 & UDP & 网络管理 \\
HTTPS & 443 & TCP & 安全Web服务 \\
RIP & 520 & UDP & 路由信息协议 \\
\bottomrule
\end{tabular}
\caption{常用协议端口号}
\end{table}

\section{协议封装层次}

\begin{table}[htbp]
\centering
\begin{tabular}{lll}
\toprule
\textbf{协议} & \textbf{封装在} & \textbf{协议号/端口} \\
\midrule
HTTP & TCP & 端口80 \\
FTP & TCP & 端口20/21 \\
SMTP & TCP & 端口25 \\
DNS & UDP & 端口53 \\
DHCP & UDP & 端口67/68 \\
RIP & UDP & 端口520 \\
TCP & IP & 协议号6 \\
UDP & IP & 协议号17 \\
ICMP & IP & 协议号1 \\
OSPF & IP & 协议号89 \\
IP & 数据链路层 & 类型0x0800 \\
ARP & 数据链路层 & 类型0x0806 \\
\bottomrule
\end{tabular}
\caption{协议封装关系}
\end{table}

\section{网络设备工作层次}

\begin{table}[htbp]
\centering
\begin{tabular}{llll}
\toprule
\textbf{设备} & \textbf{工作层次} & \textbf{隔离冲突域} & \textbf{隔离广播域} \\
\midrule
中继器 & 物理层 & 否 & 否 \\
集线器 & 物理层 & 否 & 否 \\
网桥 & 数据链路层 & 是 & 否 \\
交换机 & 数据链路层 & 是 & 否 \\
路由器 & 网络层 & 是 & 是 \\
三层交换机 & 网络层 & 是 & 是 \\
网关 & 应用层 & 是 & 是 \\
\bottomrule
\end{tabular}
\caption{网络设备对比}
\end{table}

\section{IP地址分类}

\begin{table}[htbp]
\centering\small
\begin{tabular}{ccccc}
\toprule
\textbf{类别} & \textbf{首位} & \textbf{地址范围} & \textbf{默认掩码} & \textbf{私有地址} \\
\midrule
A & 0 & 1-126 & 255.0.0.0 & 10.0.0.0/8 \\
B & 10 & 128-191 & 255.255.0.0 & 172.16.0.0/12 \\
C & 110 & 192-223 & 255.255.255.0 & 192.168.0.0/16 \\
D & 1110 & 224-239 & \multicolumn{2}{c}{多播地址} \\
E & 1111 & 240-255 & \multicolumn{2}{c}{保留地址} \\
\bottomrule
\end{tabular}
\caption{IP地址分类}
\end{table}

\section{路由协议对比}

\begin{table}[htbp]
\centering\small
\begin{tabular}{lccc}
\toprule
\textbf{特性} & \textbf{RIP} & \textbf{OSPF} & \textbf{BGP} \\
\midrule
类型 & 距离矢量 & 链路状态 & 路径矢量 \\
度量标准 & 跳数 & 链路代价 & 策略 \\
最大跳数 & 15 & 无限制 & 无限制 \\
收敛速度 & 慢 & 快 & 慢 \\
适用范围 & 小型网络 & 大型网络 & AS之间 \\
更新方式 & 周期性 & 触发式 & 触发式 \\
管理距离 & 120 & 110 & 20(eBGP)/200(iBGP) \\
封装协议 & UDP(520) & IP(89) & TCP(179) \\
\bottomrule
\end{tabular}
\caption{路由协议对比}
\end{table}

\section{TCP与UDP对比}

\begin{table}[htbp]
\centering
\begin{tabular}{lcc}
\toprule
\textbf{特性} & \textbf{TCP} & \textbf{UDP} \\
\midrule
连接性 & 面向连接 & 无连接 \\
可靠性 & 可靠传输 & 不可靠传输 \\
流量控制 & 有（滑动窗口） & 无 \\
拥塞控制 & 有 & 无 \\
首部开销 & 20字节 & 8字节 \\
传输方式 & 字节流 & 数据报 \\
适用场景 & 文件传输、邮件 & 视频、语音、DNS \\
\bottomrule
\end{tabular}
\caption{TCP与UDP对比}
\end{table}

\section{常见考点速记}

\subsection{物理层}
\begin{itemize}[leftmargin=2em,itemsep=0pt]
    \item 曼彻斯特编码：波特率是数据率的2倍
    \item 差分曼彻斯特编码：码元开始处有跳变表示0，无跳变表示1
    \item CDMA：码片序列正交，规格化内积为0
    \item 双绞线：10/100Mbps只用1、2、3、6号线
\end{itemize}

\subsection{数据链路层}
\begin{itemize}[leftmargin=2em,itemsep=0pt]
    \item 集线器：物理层，共享冲突域，不能隔离冲突域
    \item 交换机：数据链路层，隔离冲突域，不隔离广播域
    \item VLAN：隔离广播域，不同VLAN间需三层设备通信
    \item STP：防止环路，根桥优先级最小
    \item PPP：点对点协议，支持多种网络层协议
\end{itemize}

\subsection{网络层}
\begin{itemize}[leftmargin=2em,itemsep=0pt]
    \item 路由器：网络层，隔离广播域
    \item IP分片：只有源主机和路由器可分片，IPv6只有源主机分片
    \item TTL：每经过一个路由器减1，减到0则丢弃
    \item ICMP：封装在IP中，ping使用回送请求/回答
    \item ARP：将IP地址解析为MAC地址
    \item 最长前缀匹配：选择匹配位数最多的路由
\end{itemize}

\subsection{传输层}
\begin{itemize}[leftmargin=2em,itemsep=0pt]
    \item TCP三次握手：SYN→SYN+ACK→ACK
    \item TCP四次挥手：FIN→ACK→FIN→ACK
    \item TCP序号：字节序号，不是报文段序号
    \item TCP确认号：期望收到的下一字节序号
    \item 慢开始：指数增长，拥塞避免：线性增长
    \item 快重传：收到3个重复ACK立即重传
\end{itemize}

\subsection{应用层}
\begin{itemize}[leftmargin=2em,itemsep=0pt]
    \item HTTP：无状态协议，默认端口80
    \item FTP：双连接，控制连接21，数据连接20
    \item DNS：递归查询和迭代查询
    \item DHCP：动态分配IP地址，使用UDP
    \item 电子邮件：SMTP发送，POP3/IMAP接收
\end{itemize}

\section{易错点提醒}

\begin{notebox}
\textbf{单位换算}：
\begin{itemize}[itemsep=0pt]
    \item 数据量：1KB = 1024B（$2^{10}$进制）
    \item 速率：1kbps = 1000bps（$10^3$进制）
    \item 1Byte = 8bit
\end{itemize}
\end{notebox}

\begin{notebox}
\textbf{时延计算}：
\begin{itemize}[itemsep=0pt]
    \item 发送时延：与数据量和速率有关，与距离无关
    \item 传播时延：与距离和传播速度有关，与数据量无关
    \item 光纤传播速度：约$2\times10^8$m/s
    \item 电磁波传播速度：约$3\times10^8$m/s
\end{itemize}
\end{notebox}

\begin{notebox}
\textbf{子网划分}：
\begin{itemize}[itemsep=0pt]
    \item 网络地址：主机号全0
    \item 广播地址：主机号全1
    \item 可用主机数：$2^n - 2$（减去网络地址和广播地址）
    \item 子网掩码：网络位全1，主机位全0
\end{itemize}
\end{notebox}

\begin{notebox}
\textbf{协议特点}：
\begin{itemize}[itemsep=0pt]
    \item IP：无连接、不可靠、尽最大努力交付
    \item TCP：面向连接、可靠、字节流、全双工
    \item UDP：无连接、不可靠、数据报、开销小
    \item HTTP：无状态、基于TCP
\end{itemize}
\end{notebox}

\section{计算题常见类型}

\subsection{时延计算}
\begin{enumerate}[leftmargin=2em,itemsep=0pt]
    \item 计算发送时延：数据量÷速率
    \item 计算传播时延：距离÷传播速度
    \item 计算总时延：发送+传播+处理+排队
    \item 计算往返时间RTT：2×传播时延
\end{enumerate}

\subsection{信道容量计算}
\begin{enumerate}[leftmargin=2em,itemsep=0pt]
    \item 奈氏准则：$B = 2W$，再乘以每码元比特数
    \item 香农公式：先将dB转换为信噪比，再代入公式
    \item 取两者中的较小值作为实际信道容量
\end{enumerate}

\subsection{子网划分计算}
\begin{enumerate}[leftmargin=2em,itemsep=0pt]
    \item 确定需要的子网数，计算子网位数
    \item 确定需要的主机数，计算主机位数
    \item 计算子网掩码
    \item 列出各子网的地址范围
\end{enumerate}

\subsection{IP分片计算}
\begin{enumerate}[leftmargin=2em,itemsep=0pt]
    \item 计算数据部分长度：总长度-首部长度
    \item 计算每个分片最大数据：MTU-首部长度
    \item 计算分片数量：向上取整
    \item 计算片偏移：数据偏移量÷8
\end{enumerate}

\subsection{TCP序号计算}
\begin{enumerate}[leftmargin=2em,itemsep=0pt]
    \item 序号：当前序号+已发送数据字节数
    \item 确认号：期望收到的下一字节序号
    \item 三次握手：注意SYN占用一个序号
    \item 四次挥手：注意FIN占用一个序号
\end{enumerate}

\section{本章小结}

本附录汇总了计算机网络课程中的重要公式、常见端口号、协议对比、常见考点和易错点。建议在复习时：

\begin{enumerate}[leftmargin=2em,itemsep=0pt]
    \item 熟记所有公式，理解其物理意义
    \item 掌握常用协议的端口号和封装关系
    \item 理解各层协议的工作原理和特点
    \item 多做计算题，熟练掌握计算方法
    \item 注意单位换算，避免低级错误
    \item 理解概念之间的联系和区别
\end{enumerate}
